\documentclass[10pt,a4paper]{article}
	
	\usepackage{amsmath,amssymb,amsthm,mathtools,amsfonts}
	\usepackage{breqn}
	\usepackage[utf8]{inputenc}
	\usepackage[english]{babel}
	\usepackage[T1]{fontenc}
	\usepackage{graphicx}
	\graphicspath{{./img/}}
	
	%\usepackage{ebgaramond}
	\usepackage{fontspec}
	\setmainfont[Numbers=OldStyle]{EB Garamond}
	\newfontfamily\mysectionfont[Numbers=OldStyle]{Plantin MT Pro Bold Condensed}
	
	\usepackage[pdfborder={0 0 0}]{hyperref} %% Ingen firkant rundt om referencer
	\usepackage{multicol}
	\usepackage{tabularx}
	\usepackage{enumitem}% better controls with enumerating
	\usepackage{wrapfig}
		%\setlength{\wrapoverhang}{\marginparwidth}
		%\addtolength{\wrapoverhang}{\marginparsep}
		\setlength{\columnsep}{0pt}
		\setlength{\intextsep}{-10pt}
	\usepackage{caption}
	\usepackage[svgnames]{xcolor}
	\usepackage{dirtytalk}
	\usepackage[modulo]{lineno}
	\usepackage{hyperref}
	\usepackage{geometry}
	\usepackage{tasks}
	\newcommand\svar[1]{\newline\textcolor{red}{\textbf{Svar}\\#1}}
	\newcommand{\qm}[1]{‘‘#1”}
	
	\definecolor{hgblue}{RGB}{10, 40, 80}
	\definecolor{hgred}{RGB}{140, 10, 10}
	\definecolor{hggrey}{RGB}{215, 215, 210}
	\definecolor{hggreen}{RGB}{145, 205, 205}
	
	\usepackage{titlesec}
	\titleformat{\subsection}{\color{hgred}\mysectionfont\large}{}{}{}
	\newcommand{\source}[1]{Source: \href{#1}{#1}}
	\usepackage{lipsum}
	\usepackage{comment}
	\usepackage{ulem}
	
	\usepackage{bigfoot}
		\DeclareNewFootnote{a}
			\renewcommand{\thefootnotea}{\fnsymbol{footnotea}}
				\newcommand{\footnotes}[1]{\footnotea[2]{#1}} % here you can specify what symbol you want in footnotes
				\newcommand{\footnoted}[1]{\footnotea[3]{#1}} % for a second footnote on a page
				\MakePerPage{footnotes}
	
	%\reversemarginpar
	\newcommand{\glos}[3]{\dotuline{#1}\marginpar{\scriptsize\textit{#3} #2}}

	%\setlength{\parskip}{0.5\baselineskip}
	%\setlength{\parindent}{0cm}
	\setlist[enumerate]{label=\alph*),left=20pt}
	%\setlist[enumerate]{leftmargin=\parindent}
	%\setlist[enumerate]{nosep}
	
	%Titling
	\usepackage{titling}
	\setlength{\droptitle}{-12ex}
	\pretitle{\begin{flushleft}\mysectionfont\huge\color{hgblue}}
	\posttitle{\par\end{flushleft}}
	\preauthor{\begin{flushleft}\Large}
	\postauthor{\end{flushleft}}
	\predate{\begin{flushleft}}
	\postdate{\end{flushleft}}
	
	\setcounter{secnumdepth}{0}
	
	\usepackage{tikz}

\begin{document}
%\pagecolor{FloralWhite}
\title{Tommy Baton}
%\date{\vspace{-3em}} % I tilfældet ingen dato
\date{2014}
\author{Joe Goldman}
\maketitle
    
\thispagestyle{plain}
\setcounter{page}{1}
\linenumbers
\noindent"It just ain't right." Tom Odette was looking at his son Tommy practicing the baton out in the backyard. This never made him happy.


"A boy his age should be playing ball, not twirling some sissy baton." The reddish-brown hair on the back of his neck was standing on end.

Marilyn Odette was seldom at a loss for words, but this time she was torn. Tom had been trying to get her to end their son's baton-twirling days, but she just couldn't bring herself to do it. She didn't want to be the bad guy in Tommy's eyes.

"He's only eleven for godsakes." She thought at eleven years old a kid should be able to play any way he wanted as long as he wasn't hurting anybody.

"When I was eleven I was playing little league ball. I can't even look the other fathers in the face anymore, Marilyn."

She hated when Tom whined. He was too big to whine.

"It's embarrassing." He turned away from the window and looked down at his size thirteen shoe.

"You'd rather he be warming the bench at some stupid baseball game. That would make him feel real good I bet." Marilyn's words were like short biting piranhas.

"My son is turning into a goddamn fairy. How can I let him go on doing that," he said raising his voice and pointing to Tommy in the yard. "What kind of father would I be?"

"All I hear's concerns about yourself. You don't even care about Tommy."

"Yeah, I've heard it all before. I'm good for nothing. I can't get a job and I'm a failure as a father. Come on and say it. I know you want to."

He sounded wounded and Marilyn knew not to press him. Tom had been unable to find work for over eight months now. When he was laid off he tried everything he could to find work, but there was nothing out there for an otherwise unskilled livestock handler. Now he just watched TV and collected \$180.00 a week in unemployment.

"Why don't you just do guy stuff with Tommy more?" She softened her tone considerably.

"I've tried all of that. He's not interested."

Marilyn knew this. She worried about Tommy. Walnut Creek, Texas, population 1518, was a fishbowl. Everyone knew everyone else's business and everyone else's business eventually came to your doorstep whether you wanted it to or not. Marilyn had overheard catty remarks from other mothers about Tommy's baton. And Tommy's poor little eyes were starting to look permanently black and blue from near constant fights with the neighborhood kids.

From the flower-box kitchen window, she could see Tommy in the backyard practicing his spins and throwing the baton, each time higher and higher. It was a hot and still Sunday. In the glare of the harsh midday sun, Tommy looked like he was floating on air, a tiny cherub in flight. His spotless white jeans and t-shirt added to the effect. He was small for eleven. The baton was almost as long as him. She wondered when he would start to grow, and thought that he wasn't cut out for little league any more than a rooster could lay eggs. On the other hand, the beatings he took weren't worth it either.

Tommy threw the baton in the air. He began to spin, while Marilyn watched and waited for it to come down. She counted to herself. One one-thousand, two one-thousand . . .'

Tom was rustling through the refrigerator noisily as he always did when he was upset. "Hey did you buy any mustard at Kroger's yesterday? I'm gonna' make a sandwich."

"Just a minute." He was interrupting her count.

The baton came down. Tommy caught it without missing a beat and continued to twirl. Marilyn smiled and a twinkle came to her eye.

"Hot sauce!" That one must have gone up a good twenty-five or thirty feet," she beamed proudly.

Tom looked at her incredulously. "Haven't you heard anything that I've been saying here?"

"You made him this way," he said waving an accusing finger at Marilyn. "I'm putting an end to this once and for all." Storming to the back door Tom yelled, "Tommy! Get your butt in here right now!"

Marilyn was still watching Tommy out of the kitchen window. He stopped his routine, huffed, and reluctantly walked toward the back door where his father stood looming ominously. Little beads of sweat had formed on Tommy's forehead and over his upper lip, and his straight blond hair was unkempt, but other than that, he looked fresh and clean.

Because of the vast difference in their sizes, it was hard for even Marilyn to believe that Tom and Tommy were cut from the same mold.

Tom led Tommy over to the speckled formica-topped kitchen table and sat him down, pushing aside his unmade sandwich for the moment.

"Tommy boy, I think it's about time you gave up the baton. You know there's a little league try out tomorrow?"

"What, why?" Tommy was shaken and confused.

"Don't you care about all the ribbing you take over that thing? Not to mention that shiner you got there."

"Mom says people are just jealous cause I'm so good at baton."

Tom shot a cold glance at Marilyn as she busied herself with some dishes in the sink.

"But it's not something a boy your age should be good at."

"Who says?" Tommy asked a bit precociously.

Marilyn was chuckling to herself as she watched Tom fumble around for his next tact.

"Well, what about football or baseball? Don't you have any interest in sports?"

"Baton's a sport. Besides, why should I play something I'm not going to be good at?"

"Stop answering me with questions. You're getting a smart mouth; you know that? You know who says you should play boys sports? I say so. I'm taking that baton and that's final."

Tom grabbed the baton out of Tommy's hands and started into the living room.

"But I'm supposed to lead the Walnut Festival parade next month. I'm going to be on TV," Tommy shouted pleadingly.

"The hell you are. No son of mine is going on TV twirling a baton. Tomorrow we're signing you up for little league."

"Mom," Tommy begged Marilyn to intervene.

"Don't cross me, Marilyn," Tom shouted from the living room.

"I'm sorry squirt. I'm going to have to back your daddy up on this one," she said half-heartedly.

"Why Mommy, why? What did I do wrong?" Tears welled up in Tommy's eyes.

Because I got to live with him,' Marilyn thought to herself.

Marilyn couldn't stand to see Tommy cry. She remembered how happy he was the day she came home with the baton she bought at Kmart for \$7.99. It might as well have been made of solid gold. When he was younger, and Tom had the games on TV, Tommy wouldn't pay much attention until half time. Then you couldn't keep his eyes off the set. He loved to watch the baton twirlers lead the marching bands and longed to be like them.

"Nothing Tommy," Marilyn said patting him on the shoulder. "You didn't do nothing wrong. You're the best. Some things just aren't meant to be."

Tommy ran out of the room choking back his tears. He dashed through the living room and up the stairs to his room, slamming the door behind him loudly.

"Hey mister! Don't be slamming any doors in this house."

"Lighten up Tom." Marilyn snapped, walking into the living room. Making her way to the worn, floral patterned sofa, she bent down to pick up some newspapers that had fallen to the floor. She felt the effects of middle age begin to nag at her back as she straightened up slowly.

"What if we just let Tommy practice here in the backyard? He doesn't have to march in the parade."

Tom was already in his chair and hypnotized by the Cowboys game. "It's over, Mare. Decision stands."

Marilyn caught a glimpse of herself in the small Budweiser mirror that hung next to the kitchen entry as she walked listlessly out of the living room. The gray hair was starting to come in again, but she hadn't been able to afford a touch up.

In high school she had been a dark-haired beauty and Tom was the handsomest boy she'd ever seen. She used to dream that one day he would become a big Rodeo star and they would travel around the country together. Somewhere along the line she stopped thinking about that. Maybe it was the day Tom shattered his hip and doctors said his rodeo career was over. Or maybe it was when she started waiting on tables at the Fish Shanty on 1210 to help pay the bills.

She wanted things to be different for Tommy. Maybe she had indulged him a little, but he was her only child and so frail, like a sick little bird who had fallen from the nest. "You made him this way," she said mimicking Tom's accusations.

She thought about living in a place where an eleven-year-old boy wouldn't be made fun of because he twirled a baton. As hard as she tried, she just couldn't picture where that place would be.

"No there's nothing wrong with you, Tommy," she muttered to herself.

A tear welled up in the corner of her eye, but she could no longer bring herself to cry.







































\end{document}